\section{Nuclear reaction theory reconstructed from scattering principles}
\label{sec:nuclear-reaction-foundations}
% BEGIN CURATED SUBJECT INDEX
\index{optical potential}
\index{distorted-wave approximation}
\index{nuclear breakup}
% END CURATED SUBJECT INDEX

The usual vocabulary of nuclear reactions---optical potential, channel,
form factor, distorted wave, compound state, breakup---can appear to be a
collection of phenomenological devices.  It is more coherent when read as a
sequence of projections and boundary conditions.  This section derives that
sequence before any continuum discretization.  The resulting formulas are
standard; their placement in the book is not.  Each will be tied to one of the
analytic or operator-theoretic structures already constructed.

\subsection{Translational invariance and Jacobi kinetic energy}
\label{subsec:nuclear-jacobi-foundations}
% BEGIN CURATED SUBJECT INDEX
\index{central decomposition}
\index{Jacobi coordinate}
% END CURATED SUBJECT INDEX

For particles with masses $m_i$, positions $\bm r_i$, and momenta $\bm p_i$,
the laboratory kinetic energy is
\begin{equation}
  T=\sum_{i=1}^{A}\frac{\bm p_i^2}{2m_i}.
  \label{eq:many-body-lab-kinetic}
\end{equation}
Introduce the total mass, center of mass, and total momentum,
\begin{equation}
  M=\sum_i m_i,
  \qquad
  \vsub{\bm R}{cm}=\frac{1}{M}\sum_i m_i\bm r_i,
  \qquad
  \bm P=\sum_i\bm p_i.
  \label{eq:many-body-cm}
\end{equation}
A linear canonical transformation to mass-scaled Jacobi coordinates gives
\begin{equation}
  T=\frac{\bm P^2}{2M}+
  \sum_{a=1}^{A-1}\frac{\bm\pi_a^2}{2\mu_a}.
  \label{eq:jacobi-kinetic-decomposition}
\end{equation}
Removing $\vsub{\bm R}{cm}$ is not optional in a basis calculation.  If a
single-particle basis mixes intrinsic and center-of-mass excitations, spurious
levels may be misidentified as low-energy channels.

For three particles, one Jacobi set is
\begin{equation}
  \bm x=\bm r_2-\bm r_1,
  \qquad
  \bm y=\bm r_3-
  \frac{m_1\bm r_1+m_2\bm r_2}{m_1+m_2},
  \label{eq:general-three-body-jacobi}
\end{equation}
with
\begin{equation}
  \mu_x=\frac{m_1m_2}{m_1+m_2},
  \qquad
  \mu_y=\frac{m_3(m_1+m_2)}{M}.
  \label{eq:general-three-body-reduced-masses}
\end{equation}
The other two sets are obtained by cyclic permutation.  No set is physically
preferred, but a particular cluster asymptotic region is simple in one of
them.  A many-set expansion is therefore a coordinate atlas for correlations,
not an introduction of new particles.

\subsection{Cluster channels and antisymmetry}
\label{subsec:cluster-channels-pauli}
% BEGIN CURATED SUBJECT INDEX
\index{spectrum!point}
\index{scattering!Coulomb}
\index{antisymmetrization}
% END CURATED SUBJECT INDEX

Suppose the asymptotic partition is $a+A$.  A channel state has internal
states $\phi_a^{I_a}$, $\phi_A^{I_A}$ and relative orbital momentum $L$,
coupled to total $JM$:
\begin{equation}
  \mathcal Y_c^{JM}
  =\mathcal A
  \left[
   [\phi_a^{I_a}\otimes\phi_A^{I_A}]_I
   \otimes Y_L(\hat R)
  \right]_{JM},
  \qquad c=(a,A,I_a,I_A,I,L,\ldots).
  \label{eq:nuclear-channel-state}
\end{equation}
The intercluster antisymmetrizer $\mathcal A$ makes channel states nonorthogonal
at small $R$.  In a microscopic resonating-group formulation, projection
therefore gives a norm kernel as well as a Hamiltonian kernel.  In a cluster
model one often replaces this exact exchange by a Pauli projector
\begin{equation}
  \vsub{V}{Pauli}=\lambda_P
  \sum_f\ket{u_f}\bra{u_f},
  \qquad \lambda_P\to+\infty,
  \label{eq:orthogonality-condition-model}
\end{equation}
which removes forbidden relative orbits $u_f$.  The limit changes the operator
domain on the allowed subspace.  It must be held fixed while continuum model
spaces are compared.

A spectroscopic factor is a norm of a particular projected overlap and hence
depends on the resolution, interaction, and channel convention.  An
asymptotic normalization coefficient is more directly tied to the tail fixed
by a separation energy and Coulomb equation.  Reaction observables probe
these quantities through kernels; neither is an eigenvalue observable of the
full many-body Hamiltonian.

\subsection{The optical model as a Feshbach boundary value}
\label{subsec:optical-model-feshbach}
% BEGIN CURATED SUBJECT INDEX
\index{Feshbach projection}
\index{nuclear breakup}
% END CURATED SUBJECT INDEX

Let $P$ retain elastic relative motion with the target in its ground state and
let $Q$ contain every explicit excitation, transfer, breakup, and compound
configuration not retained.  Stage I gave the exact operator
\begin{equation}
  \vsub{H}{opt}(E)
  =PHP+PHQ\frac{1}{E-QHQ+\rmi0}QHP.
  \label{eq:exact-optical-H}
\end{equation}
In coordinate representation its second term is generally nonlocal,
energy dependent, angular-momentum dependent, and complex above open
thresholds.  A local Woods--Saxon or folded $g$-matrix potential is an
approximation to this object, not a fundamental interaction.

Writing $U=V+\rmi W$ and using the stationary Schr\"odinger equation gives the
continuity relation
\begin{equation}
  \nabla\cdot\bm j
  =\frac{2}{\hbar}W(\bm R)|\psi(\bm R)|^2.
  \label{eq:optical-continuity}
\end{equation}
An absorptive convention has $W\leq0$.  The lost elastic flux is the flux into
the eliminated $Q$ space only if the optical model consistently represents
that space.  When breakup channels are later included explicitly, leaving the
same phenomenological absorption untouched may count part of them twice.

For spinless elastic scattering,
\begin{equation}
  S_L=\eta_L\rme^{2\rmi\delta_L},
  \qquad 0\leq\eta_L\leq1,
  \label{eq:optical-SL}
\end{equation}
and
\begin{align}
  \vsub{\sigma}{el}
  &=\frac{\pi}{k^2}\sum_L(2L+1)|1-S_L|^2,
  \label{eq:optical-elastic-xs}\\
  \sigma_R
  &=\frac{\pi}{k^2}\sum_L(2L+1)(1-|S_L|^2).
  \label{eq:optical-reaction-xs}
\end{align}
The first measures coherent deflection; the second measures missing open-
channel flux.  A resonance can be present in either through the analytic
continuation of $S_L(E)$, but a peak in $\sigma_R$ is not by itself a pole
definition.

\subsection{Coupled channels from an internal basis}
\label{subsec:ordinary-coupled-channels}
% BEGIN CURATED SUBJECT INDEX
\index{wave packet}
\index{coupled-channel equation}
\index{CDCC}
% END CURATED SUBJECT INDEX

Let $h_A\phi_n=\epsilon_n\phi_n$ describe discrete target states.  With
\begin{equation}
  H=K_R+h_A+U(\bm R,\xi),
  \label{eq:discrete-cc-H}
\end{equation}
expand
\begin{equation}
  \Psi^{JM}
  =\frac1R\sum_{nIL}\chi_{nIL}^J(R)
  [\phi_{nI}(\xi)\otimes Y_L(\hat R)]_{JM}.
  \label{eq:discrete-cc-expansion}
\end{equation}
Projection produces
\begin{align}
  &\left[-\frac{\hbar^2}{2\mu}
   \left(\frac{\rmd^2}{\rmd R^2}-\frac{L_c(L_c+1)}{R^2}\right)
   +\epsilon_c-E\right]\chi_c^J(R)
  \notag\\
  &\qquad+\sum_{c'}U_{cc'}^J(R)\chi_{c'}^J(R)=0.
  \label{eq:ordinary-coupled-equations}
\end{align}
Thus the CDCC equation is not a special new dynamical law.  It is the ordinary
coupled-channel equation after continuum wave packets have been admitted as
internal states.

Multipole expansion gives
\begin{equation}
  U(\bm R,\xi)
  =\sum_{\lambda\mu}F_{\lambda\mu}(R,\xi)
  Y_{\lambda\mu}(\hat R),
  \label{eq:cc-multipole}
\end{equation}
and reduced matrix elements of $F_\lambda$ become radial coupling form
factors.  Angular momentum selection rules can make the coupling matrix
sparse even when the number of channels is large.  A good implementation
tests Hermitian symmetry of $U_{cc'}$ before optical absorption is introduced.

\subsection{Distorted waves and first-order transition theory}
\label{subsec:dwba-foundations}
% BEGIN CURATED SUBJECT INDEX
\index{Born approximation}
\index{boundary condition!incoming and outgoing}
\index{distorted-wave approximation}
% END CURATED SUBJECT INDEX

Choose an entrance channel Hamiltonian $H_i=K+U_i+h_i$ and write
$H=H_i+V_i$.  The exact transition amplitude to an outgoing solution of
$H_f$ is
\begin{equation}
  \vsup{T_{fi}}{prior}
  =\bra{\Psi_f^{(-)}}V_i\ket{\Phi_i\chi_i^{(+)}}.
  \label{eq:exact-prior-amplitude}
\end{equation}
Alternatively,
\begin{equation}
  \vsup{T_{fi}}{post}
  =\bra{\Phi_f\chi_f^{(-)}}V_f\ket{\Psi_i^{(+)}}.
  \label{eq:exact-post-amplitude}
\end{equation}
The two forms agree for exact wave functions with consistent interactions.
Distorted-wave Born approximation (DWBA) replaces the complete scattering
state on the right by the channel product:
\begin{equation}
  \vsup{T_{fi}}{DWBA}
  =\bra{\Phi_f\chi_f^{(-)}}\vsub{V}{tr}\ket{\Phi_i\chi_i^{(+)}}.
  \label{eq:dwba-amplitude}
\end{equation}
Post--prior disagreement then becomes a diagnostic of inconsistent optical
potentials, remnant terms, nonorthogonality, or truncation.

The distorted waves resum elastic interaction with the target.  The Born
order refers only to the transition operator $\vsub{V}{tr}$, not to replacing
all waves by planes.  In a knockout reaction, an impulse approximation further
replaces the short collision by an embedded two-nucleon $t$ matrix.  This
separation is valid when the elementary collision is short compared with the
scale over which the nuclear overlap and distortions vary; it is a scale
statement to be checked, not a definition.

\subsection{The R-matrix split: an interior connection problem}
\label{subsec:r-matrix-foundations}
% BEGIN CURATED SUBJECT INDEX
\index{boundary condition!incoming and outgoing}
\index{scattering!Coulomb}
\index{exact WKB}
\index{Feshbach projection}
\index{R-matrix@$R$-matrix}
% END CURATED SUBJECT INDEX

Choose a channel radius $a$ outside the dominant short-range interaction.
Inside, expand the wave function in square-integrable functions $X_\lambda$.
Outside, use known Coulomb or free channel functions.  A Bloch operator
\begin{equation}
  \mathcal L_B=
  \frac{\hbar^2}{2\mu}\delta(r-a)
  \left(\frac{\rmd}{\rmd r}-\frac{B}{a}\right)
  \label{eq:bloch-operator}
\end{equation}
makes the interior problem Hermitian with the chosen boundary parameter $B$.
The channel R matrix has the pole expansion
\begin{equation}
  R_{cc'}(E)
  =\sum_\lambda
  \frac{\gamma_{\lambda c}\gamma_{\lambda c'}}
       {E_\lambda-E},
  \label{eq:r-matrix-pole-expansion}
\end{equation}
where $\gamma_{\lambda c}$ is a boundary amplitude at $a$
\cite{LaneThomas1958RMatrix}.

The $E_\lambda$ depend on $a$ and $B$; observed scattering poles do not.  The
external logarithmic derivative converts the interior R matrix into the
physical collision matrix.  This is another global connection problem:
interior basis coefficients are representation data, while the matched
incoming coefficient is invariant.  Exact WKB replaces the arbitrary
interior basis by resummed local analytic solutions; the algebra of matching
is the common element.

For one isolated level and one channel, the denominator becomes
\begin{equation}
  D(E)=E_\lambda-E-\gamma_\lambda^2[S(E)-B]
       -\rmi\gamma_\lambda^2P(E),
  \label{eq:r-matrix-level-denominator}
\end{equation}
where $S$ and $P$ are shift and penetration functions.  The physical pole
solves $D(\resonantE)=0$, not simply $E=E_\lambda$.  Differentiating $D$ gives the
same energy-dependent normalization encountered for Feshbach operators.

\subsection{Three-body equations before discretization}
\label{subsec:faddeev-foundations}
% BEGIN CURATED SUBJECT INDEX
\index{Lippmann--Schwinger equation}
\index{scattering!on-shell}
\index{Faddeev equation}
\index{CDCC}
\index{nuclear breakup}
% END CURATED SUBJECT INDEX

For three particles with pair interactions $v_1,v_2,v_3$, write
\begin{equation}
  H=H_0+v_1+v_2+v_3.
  \label{eq:three-body-H-faddeev}
\end{equation}
The direct Lippmann--Schwinger kernel contains disconnected iterations and is
not compact in the useful sense.  Decompose
\begin{equation}
  \ket{\Psi}=\ket{\psi_1}+\ket{\psi_2}+\ket{\psi_3},
  \qquad
  \ket{\psi_\alpha}=G_0v_\alpha\ket{\Psi},
  \label{eq:faddeev-components}
\end{equation}
with $G_0=(E-H_0+\rmi0)^{-1}$.  Using the pair $t$ matrix
\begin{equation}
  t_\alpha=v_\alpha+v_\alpha G_0t_\alpha
  \label{eq:pair-t-matrix}
\end{equation}
gives the Faddeev equations
\begin{equation}
  \ket{\psi_\alpha}
  =G_0t_\alpha
  \sum_{\beta\neq\alpha}\ket{\psi_\beta}.
  \label{eq:faddeev-equations}
\end{equation}
Each component carries the asymptotic behavior of one pair partition, and the
sum reconstructs the common solution \cite{Faddeev1961,FaddeevMerkuriev:1993}.

The operator form for transition amplitudes is the Alt--Grassberger--Sandhas
system
\begin{equation}
  U_{\beta\alpha}
  =(1-\delta_{\beta\alpha})G_0^{-1}
  +\sum_{\gamma\neq\beta}t_\gamma G_0U_{\gamma\alpha}.
  \label{eq:ags-equations}
\end{equation}
Its on-shell matrix elements generate elastic, rearrangement, and breakup
amplitudes \cite{AltGrassbergerSandhas1967}.  CDCC does not replace the exact
three-body theory as an identity.  It approximates selected projectile
continuum and reaction couplings in a basis optimized for high-energy or
weakly bound projectile scattering.  Comparisons with Faddeev/AGS calculations
are valuable where both are feasible.

\subsection{Why three-body Coulomb is structurally harder}
\label{subsec:nuclear-three-body-coulomb}
% BEGIN CURATED SUBJECT INDEX
\index{resolvent}
\index{scattering!Coulomb}
\index{scattering!three-body Coulomb}
\index{infrared problem}
\index{Jacobi coordinate}
\index{Faddeev equation}
\index{CDCC}
% END CURATED SUBJECT INDEX

For short-range pair forces, the channel resolvents separate cleanly at large
Jacobi coordinates.  Coulomb interactions act in every asymptotic region and
leave a long-range tail in the Faddeev kernels.  Merkuriev splitting writes
each Coulomb interaction as a short-range piece adapted to the three-body
region plus a long-range piece included in the channel Hamiltonian.  The
splitting function is auxiliary; physical amplitudes must be independent of
it in a converged calculation \cite{FaddeevMerkuriev:1993}.

This is a precise nuclear version of the infrared lesson from Stage II.  A
$p+p+{}^7$Be continuum cannot be assigned one plane-wave asymptotic form.
Different Jacobi sectors carry different Coulomb distortions, and their
overlap requires matching.  CDCC with Coulomb functions in the
projectile--target coordinate controls one part of the asymptotics; CSLS or a
three-body Coulomb construction is needed to resolve fragment correlations in
the final state.

\subsection{Adiabatic and eikonal limits}
\label{subsec:adiabatic-eikonal}
% BEGIN CURATED SUBJECT INDEX
\index{scattering!Coulomb}
\index{eikonal approximation}
\index{adiabatic approximation}
\index{CDCC}
\index{nuclear breakup}
% END CURATED SUBJECT INDEX

Two scale separations lead to useful approximations.  If projectile internal
excitation energies are small compared with the collision energy over the
interaction time, replace
\begin{equation}
  h_P-\epsilon_0\longrightarrow0
  \label{eq:adiabatic-freezing}
\end{equation}
inside the reaction equation while retaining the projectile coordinates in
the interaction.  This adiabatic or sudden approximation resums breakup
configurations at degenerate internal energy.  Its failure is largest when
energy conservation among continuum channels or a narrow resonance is
important.

At high incident momentum $K$, write
\begin{equation}
  \Psi(\bm R,\xi)
  =\rme^{\rmi KZ}\widehat\Psi(\bm b,Z,\xi)
  \label{eq:eikonal-factorization}
\end{equation}
and neglect $\partial_Z^2\widehat\Psi$ relative to
$K\partial_Z\widehat\Psi$.  The equation becomes first order,
\begin{equation}
  \rmi\hbar v\frac{\partial}{\partial Z}\widehat\Psi
  =[h_P-\epsilon_0+U(\bm b,Z,\xi)]\widehat\Psi.
  \label{eq:dynamical-eikonal}
\end{equation}
Freezing $h_P$ further gives a path-ordered eikonal phase.  Equation
\eqref{eq:dynamical-eikonal} is a coupled-channel evolution along the beam
coordinate; it retains excitation energy and is often called dynamical
eikonal.

The approximation is checked by increasing energy, restricting scattering
angle, and comparing with partial-wave CDCC.  Coulomb deflection and the
logarithmic phase must be corrected explicitly; the straight-line plane-wave
eikonal is not uniformly valid at arbitrarily large impact parameter.

\subsection{Threshold scales in lithium and beryllium systems}
\label{subsec:li-be-scales}
% BEGIN CURATED SUBJECT INDEX
\index{pseudostate}
\index{nuclear breakup}
\index{lithium nuclei}
\index{beryllium nuclei}
% END CURATED SUBJECT INDEX

Weak binding creates large spatial scales
\begin{equation}
  \ell_B\sim\frac{\hbar}{\sqrt{2\mu S}},
  \label{eq:weak-binding-length}
\end{equation}
where $S$ is a separation energy.  Small $S$ means that low multipoles of the
target field couple efficiently to continuum states near threshold.  It also
means that a cluster treated as inert may have its own comparable breakup
scale.  The sequence
\begin{equation}
  {}^9\mathrm C
  \longleftrightarrow p+{}^8\mathrm B
  \longleftrightarrow p+p+{}^7\mathrm{Be}
  \label{eq:9c-threshold-hierarchy}
\end{equation}
is the central example in Laboratory VI-B.  A two-body model may reproduce one
tail but miss polarization of the intermediate $^8$B subsystem.

For $^6$Li, both $\alpha+d$ and $\alpha+p+n$ descriptions can be useful.  The
former is smaller and resolves deuteron-like breakup; the latter permits
explicit proton--neutron rearrangement, closed pseudostates, knockout overlaps,
and transition densities.  Model selection should be made by the observable
and energy scales, then verified by enlarging the space.  ``More bodies'' is
not automatically more accurate if their effective interactions and Pauli
constraints are less controlled.

\subsection{One dictionary for reaction approximations}
\label{subsec:nuclear-method-dictionary}
% BEGIN CURATED SUBJECT INDEX
\index{Wronskian}
\index{wave packet}
\index{boundary condition!incoming and outgoing}
\index{exact WKB}
\index{complex scaling}
\index{pseudostate}
\index{Feshbach projection}
\index{R-matrix@$R$-matrix}
\index{distorted-wave approximation}
\index{Faddeev equation}
% END CURATED SUBJECT INDEX

\begin{longtable}{@{}p{0.18\linewidth}p{0.24\linewidth}p{0.23\linewidth}p{0.21\linewidth}@{}}
\toprule
method & retained object & eliminated object & invariance or check\\
\midrule
\endfirsthead
\toprule
method & retained object & eliminated object & invariance or check\\
\midrule
\endhead
optical model & elastic channel & excitation and reaction space & elastic
$S$ matrix and flux deficit\\
\addlinespace
DWBA & entrance and exit distorted waves & repeated transition coupling &
post--prior and potential dependence\\
\addlinespace
R matrix & finite interior basis & exact interior continuum & channel-radius
and boundary-parameter independence\\
\addlinespace
Faddeev/\allowbreak AGS & all pair partitions & disconnected overcounting & permutation,
unitarity, and rearrangement consistency\\
\addlinespace
adiabatic and eikonal & slow internal or transverse coordinates & selected time or
longitudinal derivatives & energy and angle systematics\\
\addlinespace
CDCC & finite continuum wave packets & $Q_N$ continuum & model-space and
smoothing convergence\\
\addlinespace
complex scaling & exposed pole plus rotated cut & undeformed outgoing growth &
angle and contour independence\\
\addlinespace
exact WKB & resummed connection datum & divergent local formal series & Stokes,
Wronskian, and homotopy consistency\\
\bottomrule
\end{longtable}

The table prevents a common misuse of comparisons.  Two methods can disagree
because they approximate different eliminated spaces, not because one is
simply of higher ``order.''  To compare them, match the retained Hamiltonian,
asymptotic channels, and observable first.  Only then vary the approximation
specific regulator or truncation.

\begin{checkpointbox}
The reader should now be able to derive optical absorption from a Feshbach
boundary value, obtain ordinary coupled-channel equations, state the
post--prior test of DWBA, interpret R-matrix levels as interior representation
data, and explain why Faddeev components are adapted to different asymptotic
partitions.  CDCC will next be obtained by replacing the discrete internal
basis of Eq.~\eqref{eq:discrete-cc-expansion} with normalized continuum bins or
pseudostates.
\end{checkpointbox}
